\ifdefined\inMainDoc\else
\documentclass[12pt,a4paper]{article}
% ================================================================
%  PRÉAMBULE COMMUN - ANNALES CORRIGÉES
%  Université de Kara - FaST - Licence Fondamentale MPC
%  Design optimisé impression noir et blanc
% ================================================================

% -- ENCODAGE & LANGUE --
\usepackage[utf8]{inputenc}
\usepackage[T1]{fontenc}
\usepackage[french]{babel}
\usepackage{lmodern}
\usepackage{microtype}

% -- MISE EN PAGE --
\usepackage[a4paper, top=2.8cm, bottom=2.5cm, left=2.5cm, right=2.5cm, headheight=14pt]{geometry}
\usepackage{setspace}
\setstretch{1.2}
\usepackage{parskip}
\setlength{\parindent}{1.2em}
\setlength{\parskip}{4pt}

% -- MATHS --
\usepackage{amsmath, amssymb, mathtools, amsthm}
\usepackage{mathrsfs}
\usepackage{dsfont}
\usepackage{bm}
\usepackage{cancel}
\usepackage{textcomp}
% Définition de \oiint (intégrale de surface fermée) sans package supplémentaire
\newcommand{\oiint}{\mathop{\iint\!\!\!\!\!\!\!\bigcirc}\limits}

% -- PHYSIQUE & CHIMIE --
\usepackage{physics}
\usepackage{siunitx}
% Lorsque physics est chargé, siunitx omet \qty. On le restaure ici :
\AtBeginDocument{\RenewCommandCopy\qty\SI}
\sisetup{locale = FR, detect-all, output-decimal-marker={,}}
% Commande de substitution pour les unités posant problème avec pdflatex
% (ohm, degreeCelsius, micro, angstrom, percent utilisent des caractères non-ASCII)
\newcommand{\qtyohm}[1]{\ensuremath{\num{#1}\,\Omega}}
\newcommand{\qtydegc}[1]{\ensuremath{\num{#1}\,{}^\circ\text{C}}}
\newcommand{\qtyangstrom}[1]{\ensuremath{\num{#1}\,\text{\AA}}}
\newcommand{\qtypercent}[1]{\ensuremath{\num{#1}\,\%}}
\newcommand{\qtyum}[1]{\ensuremath{\num{#1}\,\mu\text{m}}}
\usepackage[version=4]{mhchem}

% -- GRAPHISMES --
\usepackage{graphicx}
\usepackage{float}
\usepackage{caption}
\usepackage{subcaption}
\usepackage{wrapfig}
\usepackage{tikz}
\usetikzlibrary{calc, arrows.meta, patterns, shapes.geometric}
\usepackage{pgfplots}
\pgfplotsset{compat=1.18}

% -- TABLEAUX --
\usepackage{booktabs}
\usepackage{array}
\usepackage{tabularx}
\usepackage{multirow}

% -- LISTES --
\usepackage{enumitem}
\setlist[enumerate]{topsep=3pt, itemsep=2pt, parsep=0pt}
\setlist[itemize]{topsep=3pt, itemsep=2pt, parsep=0pt, label={.}}

% -- BOÎTES (noir et blanc) --
\usepackage{mdframed}
\usepackage{tcolorbox}
\tcbuselibrary{skins, breakable, theorems}

% Boîte EXERCICE : encadré avec filet épais, fond blanc
\newmdenv[
  linewidth=1.2pt,
  linecolor=black,
  roundcorner=3pt,
  skipabove=10pt,
  skipbelow=6pt,
  innertopmargin=8pt,
  innerbottommargin=8pt,
  innerleftmargin=10pt,
  innerrightmargin=10pt,
  backgroundcolor=white,
  frametitle={\bfseries\small Exercice},
  frametitlealignment=\raggedright,
  frametitleaboveskip=4pt,
  frametitlebelowskip=4pt,
]{boiteexercice}

% Boîte CORRIGÉ : fond gris très clair + filet fin
\newmdenv[
  linewidth=0.6pt,
  linecolor=black,
  leftline=true,
  rightline=false,
  topline=false,
  bottomline=false,
  linecolor=black,
  innerleftmargin=12pt,
  innerrightmargin=8pt,
  innertopmargin=6pt,
  innerbottommargin=6pt,
  backgroundcolor=gray!8,
  skipabove=4pt,
  skipbelow=8pt,
]{boitecorrige}

% Boîte RÉSUMÉ DE COURS : double encadré
\newmdenv[
  linewidth=0.8pt,
  linecolor=black,
  roundcorner=2pt,
  backgroundcolor=gray!12,
  innertopmargin=10pt,
  innerbottommargin=10pt,
  innerleftmargin=12pt,
  innerrightmargin=12pt,
  skipabove=12pt,
  skipbelow=8pt,
]{boiteresume}

% Boîte ATTENTION/REMARQUE : barre gauche épaisse
\newmdenv[
  leftline=true,
  rightline=false,
  topline=false,
  bottomline=false,
  linewidth=3pt,
  linecolor=black,
  backgroundcolor=gray!5,
  innerleftmargin=12pt,
  innerrightmargin=8pt,
  innertopmargin=5pt,
  innerbottommargin=5pt,
  skipabove=6pt,
  skipbelow=6pt,
]{boiteremarque}

% -- DIFFICULTÉ DES EXERCICES --
\newcommand{\facile}{$\star$}
\newcommand{\moyen}{$\star\!\star$}
\newcommand{\difficile}{$\star\!\star\!\star$}
\newcommand{\niveau}[1]{\hfill\textbf{#1}\hspace{2pt}}

% -- EN-TÊTES ET PIEDS DE PAGE --
\usepackage{fancyhdr}
\pagestyle{fancy}
\fancyhf{}
\renewcommand{\headrulewidth}{0.5pt}
\renewcommand{\footrulewidth}{0.3pt}

% -- HYPERLIENS --
\usepackage[hidelinks]{hyperref}
\hypersetup{
    colorlinks=false,
    pdftitle={Annale Corrigée - Licence MPC - Université de Kara}
}

% -- TITRES --
\usepackage{titlesec}
\titleformat{\section}{\Large\bfseries}{\thesection.}{0.8em}{}[\vspace{2pt}\hrule\vspace{4pt}]
\titleformat{\subsection}{\large\bfseries}{\thesubsection.}{0.7em}{}
\titleformat{\subsubsection}{\normalsize\bfseries}{\thesubsubsection.}{0.6em}{}

% -- THÉORÈMES --
\theoremstyle{definition}
\newtheorem{defn}{Définition}[section]
\newtheorem{prop}{Propriété}[section]
\newtheorem{thm}{Théorème}[section]
\newtheorem{corol}{Corollaire}[section]
\newtheorem{exemple}{Exemple}[section]

% -- COMMANDES MATHÉMATIQUES UTILES --
\newcommand{\vect}[1]{\overrightarrow{#1}}
\newcommand{\R}{\mathbb{R}}
\newcommand{\N}{\mathbb{N}}
\newcommand{\Z}{\mathbb{Z}}
\newcommand{\C}{\mathbb{C}}
\renewcommand{\d}{\mathrm{d}}
\newcommand{\deriv}[2]{\dfrac{\d #1}{\d #2}}
\newcommand{\dpart}[2]{\dfrac{\partial #1}{\partial #2}}
\newcommand{\rot}{\overrightarrow{\mathrm{rot}}\,}
\renewcommand{\div}{\mathrm{div}\,}
\renewcommand{\grad}{\overrightarrow{\mathrm{grad}}\,}
\newcommand{\e}{\mathrm{e}}
\newcommand{\ii}{\mathrm{i}}
\renewcommand{\Tr}{\mathrm{Tr}}

% Numérotation des équations par section
\numberwithin{equation}{section}

% -- RÉSULTATS ENCADRÉS --
\newcommand{\res}[1]{\[\boxed{#1}\]}

% -- REMARQUE INLINE --
\newcommand{\rmq}[1]{\begin{boiteremarque}\textbf{Remarque.} #1\end{boiteremarque}}

% -- MISE EN VALEUR DU RÉSULTAT FINAL --
\newcommand{\reponse}[1]{%
  \begin{center}
    \fbox{\begin{minipage}{0.92\textwidth}\centering #1\end{minipage}}
  \end{center}%
}


\lhead{\small\textit{Analyse Vectorielle - Licence 1}}
\rhead{\small\textit{FaST, Université de Kara}}
\cfoot{\thepage}

\begin{document}
\fi

\begin{center}
  {\LARGE\bfseries ANALYSE VECTORIELLE}\\[0.3cm]
  {\normalsize Licence 1 - Mathématiques, Physique et Chimie}\\[0.1cm]
  \rule{0.7\textwidth}{0.8pt}
\end{center}

\section*{Résumé du cours}

\begin{boiteresume}

\textbf{1. Opérateurs différentiels en coordonnées cartésiennes.} Pour un champ scalaire $f(x,y,z)$ et un champ vectoriel $\vec{A}=(A_x,A_y,A_z)$ :
\begin{itemize}[nosep]
  \item Gradient : $\overrightarrow{\text{grad}}\,f = \frac{\partial f}{\partial x}\,\vec{i} + \frac{\partial f}{\partial y}\,\vec{j} + \frac{\partial f}{\partial z}\,\vec{k}$
  \item Divergence : $\text{div}\,\vec{A} = \frac{\partial A_x}{\partial x}+\frac{\partial A_y}{\partial y}+\frac{\partial A_z}{\partial z}$
  \item Rotationnel : $\overrightarrow{\text{rot}}\,\vec{A} = \left(\frac{\partial A_z}{\partial y}-\frac{\partial A_y}{\partial z}\right)\vec{i} + \left(\frac{\partial A_x}{\partial z}-\frac{\partial A_z}{\partial x}\right)\vec{j} + \left(\frac{\partial A_y}{\partial x}-\frac{\partial A_x}{\partial y}\right)\vec{k}$
  \item Laplacien : $\Delta f = \frac{\partial^2 f}{\partial x^2}+\frac{\partial^2 f}{\partial y^2}+\frac{\partial^2 f}{\partial z^2}$
\end{itemize}

\textbf{2. Propriétés fondamentales.}
\begin{itemize}[nosep]
  \item $\overrightarrow{\text{rot}}(\overrightarrow{\text{grad}}\,f) = \vec{0}$ (le gradient est irrotationnel)
  \item $\text{div}(\overrightarrow{\text{rot}}\,\vec{A}) = 0$ (le rotationnel est solénoïdal)
  \item Si $\overrightarrow{\text{rot}}\,\vec{A}=\vec{0}$, $\vec{A}$ dérive d'un potentiel scalaire : $\vec{A}=\overrightarrow{\text{grad}}\,f$
\end{itemize}

\textbf{3. Théorème de Green-Ostrogradski (divergence).} Pour un volume $V$ de frontière $S$ :
\[
\iiint_V \text{div}\,\vec{A}\,\d V = \oiint_S \vec{A}\cdot\d\vec{S}
\]

\textbf{4. Théorème de Stokes.} Pour une surface $S$ de contour orienté $C$ :
\[
\iint_S \overrightarrow{\text{rot}}\,\vec{A}\cdot\d\vec{S} = \oint_C \vec{A}\cdot\d\vec{l}
\]

\textbf{5. Coordonnées cylindriques.} $(r,\theta,z)$ : $\overrightarrow{\text{grad}}\,f = \frac{\partial f}{\partial r}\vec{e}_r + \frac{1}{r}\frac{\partial f}{\partial\theta}\vec{e}_\theta + \frac{\partial f}{\partial z}\vec{e}_z$, $\;\text{div}\,\vec{A} = \frac{1}{r}\frac{\partial(rA_r)}{\partial r}+\frac{1}{r}\frac{\partial A_\theta}{\partial\theta}+\frac{\partial A_z}{\partial z}$.

\textbf{6. Coordonnées sphériques.} $(r,\theta,\varphi)$ : $\overrightarrow{\text{grad}}\,f = \frac{\partial f}{\partial r}\vec{e}_r + \frac{1}{r}\frac{\partial f}{\partial\theta}\vec{e}_\theta + \frac{1}{r\sin\theta}\frac{\partial f}{\partial\varphi}\vec{e}_\varphi$, $\;\text{div}\,\vec{A} = \frac{1}{r^2}\frac{\partial(r^2A_r)}{\partial r}+\frac{1}{r\sin\theta}\frac{\partial(\sin\theta\,A_\theta)}{\partial\theta}+\frac{1}{r\sin\theta}\frac{\partial A_\varphi}{\partial\varphi}$.

\end{boiteresume}

\newpage
\section{Opérateurs différentiels}

\subsection*{Exercice 1 - Calcul de gradient, divergence et rotationnel \niveau{\facile}}

\begin{boiteexercice}
Calculer le gradient, la divergence ou le rotationnel selon le cas :
\begin{enumerate}[label=(\alph*)]
  \item $\overrightarrow{\text{grad}}\,f$ avec $f(x,y,z) = x^2y + yz^2$
  \item $\text{div}\,\vec{A}$ avec $\vec{A} = (xy, y^2z, z^3)$
  \item $\overrightarrow{\text{rot}}\,\vec{B}$ avec $\vec{B} = (y, -x, 0)$
  \item Le laplacien $\Delta f$ avec $f(x,y,z) = \e^x\sin y$
\end{enumerate}
\end{boiteexercice}

\begin{boitecorrige}
\textbf{(a)}
\[
\overrightarrow{\text{grad}}\,f = \frac{\partial f}{\partial x}\vec{i}+\frac{\partial f}{\partial y}\vec{j}+\frac{\partial f}{\partial z}\vec{k} = 2xy\,\vec{i} + (x^2+z^2)\,\vec{j} + 2yz\,\vec{k}
\]

\textbf{(b)}
\[
\text{div}\,\vec{A} = \frac{\partial(xy)}{\partial x}+\frac{\partial(y^2z)}{\partial y}+\frac{\partial(z^3)}{\partial z} = y + 2yz + 3z^2
\]

\textbf{(c)}
\[
\overrightarrow{\text{rot}}\,\vec{B} = \begin{vmatrix}\vec{i}&\vec{j}&\vec{k}\\\partial_x&\partial_y&\partial_z\\y&-x&0\end{vmatrix}
= \left(\frac{\partial 0}{\partial y}-\frac{\partial(-x)}{\partial z}\right)\vec{i} - \left(\frac{\partial 0}{\partial x}-\frac{\partial y}{\partial z}\right)\vec{j} + \left(\frac{\partial(-x)}{\partial x}-\frac{\partial y}{\partial y}\right)\vec{k}
\]
\[
= (0-0)\vec{i} - (0-0)\vec{j} + (-1-1)\vec{k} = -2\,\vec{k}
\]
Ce résultat correspond à un champ de vitesse en rotation rigide : $\vec{B} = \vec{\Omega}\wedge\vec{r}$ avec $\vec{\Omega}=-\vec{k}$.

\textbf{(d)}
\[
\frac{\partial f}{\partial x}=\e^x\sin y,\quad \frac{\partial^2 f}{\partial x^2}=\e^x\sin y
\]
\[
\frac{\partial f}{\partial y}=\e^x\cos y,\quad \frac{\partial^2 f}{\partial y^2}=-\e^x\sin y
\]
\[
\frac{\partial^2 f}{\partial z^2}=0
\]
\[
\Delta f = \e^x\sin y - \e^x\sin y + 0 = 0
\]
La fonction $\e^x\sin y$ est harmonique (vérifie $\Delta f=0$).
\end{boitecorrige}

\subsection*{Exercice 2 - Champ conservatif \niveau{\moyen}}

\begin{boiteexercice}
On considère le champ $\vec{F} = (2xy+z^2,\;x^2,\;2xz)$.
\begin{enumerate}
  \item Montrer que $\vec{F}$ est irrotationnel.
  \item En déduire qu'il existe un potentiel scalaire $\varphi$ tel que $\vec{F}=\overrightarrow{\text{grad}}\,\varphi$. Déterminer $\varphi$.
  \item Calculer la circulation de $\vec{F}$ le long d'un chemin allant de $A=(0,0,0)$ à $B=(1,1,1)$.
\end{enumerate}
\end{boiteexercice}

\begin{boitecorrige}
\textbf{1.} Calculons $\overrightarrow{\text{rot}}\,\vec{F}$ :
\[
(\text{rot}\,\vec{F})_x = \frac{\partial(2xz)}{\partial y}-\frac{\partial(x^2)}{\partial z} = 0-0=0
\]
\[
(\text{rot}\,\vec{F})_y = \frac{\partial(2xy+z^2)}{\partial z}-\frac{\partial(2xz)}{\partial x} = 2z-2z=0
\]
\[
(\text{rot}\,\vec{F})_z = \frac{\partial(x^2)}{\partial x}-\frac{\partial(2xy+z^2)}{\partial y} = 2x-2x=0
\]
Le champ est irrotationnel.

\textbf{2.} On cherche $\varphi$ tel que $\partial\varphi/\partial x = 2xy+z^2$, $\partial\varphi/\partial y=x^2$, $\partial\varphi/\partial z=2xz$.

Intégration selon $x$ : $\varphi = x^2y + xz^2 + g(y,z)$.

Dérivation selon $y$ : $\partial\varphi/\partial y = x^2 + \partial g/\partial y = x^2$, donc $\partial g/\partial y=0$, soit $g=h(z)$.

Dérivation selon $z$ : $\partial\varphi/\partial z = 2xz + h'(z) = 2xz$, donc $h'(z)=0$, $h=\text{cste}$.
\[
\boxed{\varphi(x,y,z) = x^2y + xz^2 + C}
\]

\textbf{3.} Pour un champ conservatif, la circulation ne dépend que des extrémités :
\[
\int_A^B \vec{F}\cdot\d\vec{l} = \varphi(B)-\varphi(A) = \varphi(1,1,1)-\varphi(0,0,0) = (1+1)-0 = 2
\]
\end{boitecorrige}

\section{Théorèmes intégraux}

\subsection*{Exercice 3 - Théorème de Green-Ostrogradski \niveau{\moyen}}

\begin{boiteexercice}
Calculer le flux du champ $\vec{A} = (x^2, y^2, z^2)$ à travers la surface du cube $[0,1]^3$ :
\begin{enumerate}
  \item Directement (en calculant le flux sur chaque face).
  \item En appliquant le théorème de la divergence.
\end{enumerate}
\end{boiteexercice}

\begin{boitecorrige}
\textbf{1. Calcul direct.} Le cube a 6 faces. Par symétrie, les faces $x=0$, $y=0$, $z=0$ contribuent toutes avec $\vec{A}\cdot(-\vec{n})=0$ (car $A_x|_{x=0}=0$, etc.). Les faces $x=1$, $y=1$, $z=1$ donnent :
\[
\Phi_{x=1} = \int_0^1\int_0^1 1^2\,\d y\,\d z = 1, \quad \Phi_{y=1} = 1, \quad \Phi_{z=1} = 1
\]
Flux total : $\Phi = 1+1+1 = 3$.

\textbf{2. Théorème de la divergence.}
\[
\text{div}\,\vec{A} = \frac{\partial x^2}{\partial x}+\frac{\partial y^2}{\partial y}+\frac{\partial z^2}{\partial z} = 2x+2y+2z
\]
\begin{align*}
\Phi &= \iiint_{[0,1]^3}(2x+2y+2z)\,\d x\,\d y\,\d z \\
&= 3\int_0^1\int_0^1\int_0^1 2x\,\d x\,\d y\,\d z = 3\times1\times1\times1 = 3
\end{align*}
Les deux méthodes donnent bien le même résultat.
\end{boitecorrige}

\subsection*{Exercice 4 - Théorème de Stokes \niveau{\moyen}}

\begin{boiteexercice}
Soit $\vec{F} = (-y, x, 0)$ et $S$ le disque unité $x^2+y^2\leq1$, $z=0$, orienté vers le haut.
\begin{enumerate}
  \item Calculer $\overrightarrow{\text{rot}}\,\vec{F}$.
  \item Calculer $\iint_S \overrightarrow{\text{rot}}\,\vec{F}\cdot\d\vec{S}$.
  \item Retrouver ce résultat en calculant la circulation $\oint_C \vec{F}\cdot\d\vec{l}$ sur le cercle unité.
\end{enumerate}
\end{boiteexercice}

\begin{boitecorrige}
\textbf{1.}
\begin{align*}
\overrightarrow{\text{rot}}\,\vec{F} &= \begin{vmatrix}\vec{i}&\vec{j}&\vec{k}\\\partial_x&\partial_y&\partial_z\\-y&x&0\end{vmatrix}
= (0-0)\vec{i}-(0-0)\vec{j}+\left(\frac{\partial x}{\partial x}-\frac{\partial(-y)}{\partial y}\right)\vec{k} \\
&= (1+1)\vec{k} = 2\vec{k}
\end{align*}

\textbf{2.} Sur $S$, $\d\vec{S}=\d x\,\d y\,\vec{k}$, donc :
\[
\iint_S \overrightarrow{\text{rot}}\,\vec{F}\cdot\d\vec{S} = \iint_S 2\,\d x\,\d y = 2\times\text{Aire}(S) = 2\pi
\]

\textbf{3.} Sur le cercle unité, $\vec{r}(\theta)=(\cos\theta,\sin\theta,0)$, $\d\vec{l}=(-\sin\theta,\cos\theta,0)\,\d\theta$ :
\[
\oint_C \vec{F}\cdot\d\vec{l} = \int_0^{2\pi}(-\sin\theta,\cos\theta,0)\cdot(-\sin\theta,\cos\theta,0)\,\d\theta = \int_0^{2\pi}(\sin^2\theta+\cos^2\theta)\,\d\theta = 2\pi
\]
On retrouve bien $2\pi$, conformément au théorème de Stokes.
\end{boitecorrige}

\subsection*{Exercice 5 - Coordonnées cylindriques \niveau{\difficile}}

\begin{boiteexercice}
En coordonnées cylindriques $(r,\theta,z)$, soit le champ $\vec{A} = A_r(r)\,\vec{e}_r$ à symétrie cylindrique.
\begin{enumerate}
  \item Exprimer la divergence de $\vec{A}$.
  \item On suppose $\text{div}\,\vec{A} = f(r)$. Exprimer $A_r(r)$ en fonction de $f$.
  \item Application : si $\vec{A}$ représente un champ électrique créé par une distribution volumique $\rho = \rho_0$ (uniforme) dans le cylindre $r\leq R$, retrouver le champ à l'intérieur et à l'extérieur.
\end{enumerate}
\end{boiteexercice}

\begin{boitecorrige}
\textbf{1.} En coordonnées cylindriques, pour un champ purement radial ($A_\theta=A_z=0$) :
\[
\text{div}\,\vec{A} = \frac{1}{r}\frac{\partial(r A_r)}{\partial r}
\]

\textbf{2.} $\frac{1}{r}\frac{\d(rA_r)}{\d r} = f(r)$, donc $\frac{\d(rA_r)}{\d r} = rf(r)$. En intégrant :
\[
rA_r = \int_0^r r'f(r')\,\d r', \quad A_r(r) = \frac{1}{r}\int_0^r r'f(r')\,\d r'
\]

\textbf{3.} Par le théorème de Gauss, $\vec{E}=E_r(r)\vec{e}_r$ et $\text{div}\,\vec{E}=\rho/\varepsilon_0$.

\textit{Intérieur} ($r\leq R$) : $\frac{1}{r}\frac{\d(rE_r)}{\d r}=\frac{\rho_0}{\varepsilon_0}$, donc $rE_r=\frac{\rho_0 r^2}{2\varepsilon_0}+C$. Avec $E_r(0)$ fini, $C=0$ :
\[
E_r = \frac{\rho_0 r}{2\varepsilon_0}
\]

\textit{Extérieur} ($r>R$) : pas de charge. $\frac{\d(rE_r)}{\d r}=0$ donc $rE_r=\text{cste}=\frac{\rho_0 R^2}{2\varepsilon_0}$ (continuité en $r=R$) :
\[
E_r = \frac{\rho_0 R^2}{2\varepsilon_0 r}
\]
Ce résultat est l'analogue cylindrique du théorème de Gauss pour une distribution uniforme.
\end{boitecorrige}

\subsection*{Exercice 6 - Opérateurs en coordonnées sphériques \niveau{\difficile}}

\begin{boiteexercice}
En coordonnées sphériques $(r,\theta,\varphi)$, soit le champ $f(r) = \frac{1}{r}$ (potentiel coulombien, $r\neq0$).
\begin{enumerate}
  \item Calculer $\overrightarrow{\text{grad}}\,f$.
  \item Calculer $\Delta f = \text{div}(\overrightarrow{\text{grad}}\,f)$ pour $r\neq0$.
  \item Montrer que le champ $\vec{E} = -\overrightarrow{\text{grad}}\,f$ satisfait $\text{div}\,\vec{E} = 0$ pour $r\neq0$.
\end{enumerate}
\end{boiteexercice}

\begin{boitecorrige}
\textbf{1.} $f$ ne dépend que de $r$, donc en coordonnées sphériques :
\[
\overrightarrow{\text{grad}}\,f = \frac{\partial f}{\partial r}\,\vec{e}_r = -\frac{1}{r^2}\,\vec{e}_r
\]

\textbf{2.} Le laplacien d'un champ scalaire à symétrie sphérique est :
\[
\Delta f = \frac{1}{r^2}\frac{\d}{\d r}\left(r^2\frac{\d f}{\d r}\right) = \frac{1}{r^2}\frac{\d}{\d r}\left(r^2\cdot\left(-\frac{1}{r^2}\right)\right) = \frac{1}{r^2}\frac{\d}{\d r}(-1) = 0
\]
La fonction $f = 1/r$ est \textbf{harmonique} pour $r\neq0$.

\textbf{3.} $\vec{E} = -\overrightarrow{\text{grad}}\,f = \frac{1}{r^2}\,\vec{e}_r$. Sa divergence :
\[
\text{div}\,\vec{E} = \frac{1}{r^2}\frac{\d(r^2 E_r)}{\d r} = \frac{1}{r^2}\frac{\d}{\d r}\left(r^2\cdot\frac{1}{r^2}\right) = \frac{1}{r^2}\frac{\d(1)}{\d r} = 0
\]
Ce résultat est lié au fait que ce champ est créé par une charge ponctuelle à l'origine : il est sans divergence \textit{partout sauf en} $r=0$.
\end{boitecorrige}

\subsection*{Exercice 7 - Identités vectorielles \niveau{\moyen}}

\begin{boiteexercice}
Vérifier les identités vectorielles suivantes pour des champs scalaires $f$, $g$ et des champs vectoriels $\vec{A}$, $\vec{B}$ :
\begin{enumerate}[label=(\alph*)]
  \item $\overrightarrow{\text{grad}}(fg) = f\,\overrightarrow{\text{grad}}\,g + g\,\overrightarrow{\text{grad}}\,f$
  \item $\text{div}(f\vec{A}) = f\,\text{div}\,\vec{A} + \vec{A}\cdot\overrightarrow{\text{grad}}\,f$
  \item $\overrightarrow{\text{rot}}(f\vec{A}) = f\,\overrightarrow{\text{rot}}\,\vec{A} + \overrightarrow{\text{grad}}\,f \wedge \vec{A}$
  \item Application : si $\vec{E} = -\overrightarrow{\text{grad}}\,V$ et $\vec{D} = \varepsilon\vec{E}$ ($\varepsilon$ scalaire), montrer que $\text{div}\,\vec{D} = -\varepsilon\Delta V$.
\end{enumerate}
\end{boiteexercice}

\begin{boitecorrige}
\textbf{(a)} Composante selon $x$ : $\frac{\partial(fg)}{\partial x} = f\frac{\partial g}{\partial x}+g\frac{\partial f}{\partial x}$ (règle du produit). Identique pour $y$ et $z$.

\textbf{(b)} $\text{div}(f\vec{A}) = \frac{\partial(fA_x)}{\partial x}+\frac{\partial(fA_y)}{\partial y}+\frac{\partial(fA_z)}{\partial z}$. Par la règle du produit :
\[
= f\frac{\partial A_x}{\partial x}+A_x\frac{\partial f}{\partial x}+\cdots = f\,\text{div}\,\vec{A} + \vec{A}\cdot\overrightarrow{\text{grad}}\,f
\]

\textbf{(c)} Composante $z$ de $\overrightarrow{\text{rot}}(f\vec{A})$ :
\[
\frac{\partial(fA_y)}{\partial x}-\frac{\partial(fA_x)}{\partial y} = f\frac{\partial A_y}{\partial x}+A_y\frac{\partial f}{\partial x}-f\frac{\partial A_x}{\partial y}-A_x\frac{\partial f}{\partial y}
\]
\[
= f\left(\frac{\partial A_y}{\partial x}-\frac{\partial A_x}{\partial y}\right)+\left(A_y\frac{\partial f}{\partial x}-A_x\frac{\partial f}{\partial y}\right) = f(\overrightarrow{\text{rot}}\,\vec{A})_z + (\overrightarrow{\text{grad}}\,f\wedge\vec{A})_z
\]

\textbf{(d)} En appliquant (b) avec $\vec{A}=\vec{E}=-\overrightarrow{\text{grad}}\,V$ et $f=\varepsilon$ (constant) :
\[
\text{div}\,\vec{D} = \text{div}(\varepsilon\vec{E}) = \varepsilon\,\text{div}\,\vec{E}+\vec{E}\cdot\overrightarrow{\text{grad}}\,\varepsilon
\]
Si $\varepsilon$ est constant, $\overrightarrow{\text{grad}}\,\varepsilon = \vec{0}$ et $\text{div}\,\vec{E} = \text{div}(-\overrightarrow{\text{grad}}\,V) = -\Delta V$.
Donc $\text{div}\,\vec{D} = -\varepsilon\Delta V$.
\end{boitecorrige}

\subsection*{Exercice 8 - Flux à travers un cylindre \niveau{\moyen}}

\begin{boiteexercice}
Calculer le flux du champ $\vec{A} = (x, y, z)$ à travers la surface latérale du cylindre de rayon $R=2$ et de hauteur $H=3$ $(0\leq z\leq3)$, puis par le théorème de la divergence (en incluant les deux disques).
\end{boiteexercice}

\begin{boitecorrige}
\textbf{Par le théorème de la divergence :}
\[
\text{div}\,\vec{A} = \frac{\partial x}{\partial x}+\frac{\partial y}{\partial y}+\frac{\partial z}{\partial z} = 3
\]
\[
\Phi_{total} = \iiint_V 3\,\d V = 3\times\text{Volume} = 3\times\pi R^2 H = 3\pi\times4\times3 = 36\pi
\]

\textbf{Flux sur les disques :}
\begin{itemize}[nosep]
  \item Disque supérieur ($z=3$, $\vec{n}=+\hat{z}$) : $\Phi_{sup} = \iint(x,y,3)\cdot(0,0,1)\,\d S = 3\pi R^2 = 12\pi$.
  \item Disque inférieur ($z=0$, $\vec{n}=-\hat{z}$) : $\Phi_{inf} = \iint(x,y,0)\cdot(0,0,-1)\,\d S = 0$.
\end{itemize}

\textbf{Flux sur la surface latérale :}
Sur le cylindre $r=R=2$, $\vec{n} = \hat{r} = (x/R, y/R, 0)$ et $A_n = \vec{A}\cdot\hat{r} = (x^2+y^2)/R = R = 2$.
\[
\Phi_{lat} = \iint_S A_n\,\d S = 2\times(2\pi R H) = 2\times(2\pi\times2\times3) = 24\pi
\]

Vérification : $\Phi_{sup}+\Phi_{inf}+\Phi_{lat} = 12\pi+0+24\pi = 36\pi$.
\end{boitecorrige}

\subsection*{Exercice 9 - Champ solénoïdal \niveau{\moyen}}

\begin{boiteexercice}
Un champ vectoriel $\vec{B}$ est dit solénoïdal si $\text{div}\,\vec{B} = 0$.
\begin{enumerate}
  \item Vérifier que $\vec{B}_1 = (y, -x, 0)$ est solénoïdal.
  \item Vérifier que $\vec{B}_2 = (yz, xz, xy)$ est solénoïdal.
  \item Montrer que $\vec{B}_2 = \overrightarrow{\text{rot}}\,\vec{A}$ avec $\vec{A} = (0, 0, xy)$ (vérification que le rotationnel d'un champ est toujours solénoïdal).
  \item Utiliser le théorème de Stokes pour calculer $\iint_S \vec{B}_1\cdot\d\vec{S}$ où $S$ est l'hémisphère unité $x^2+y^2+z^2=1$, $z\geq0$, orientée vers le haut.
\end{enumerate}
\end{boiteexercice}

\begin{boitecorrige}
\textbf{1.} $\text{div}\,\vec{B}_1 = \frac{\partial y}{\partial x}+\frac{\partial(-x)}{\partial y}+0 = 0+0 = 0$.

\textbf{2.} $\text{div}\,\vec{B}_2 = \frac{\partial(yz)}{\partial x}+\frac{\partial(xz)}{\partial y}+\frac{\partial(xy)}{\partial z} = 0+0+0 = 0$.

\textbf{3.} $\overrightarrow{\text{rot}}(0,0,xy)$ :
\[
\text{rot}_x = \frac{\partial(xy)}{\partial y}-\frac{\partial 0}{\partial z} = x-0 = x \neq yz \text{ en général}.
\]
Prenons $\vec{A} = (0, 0, xy)$ et calculons :
\[
\overrightarrow{\text{rot}}\,\vec{A} = \left(\frac{\partial(xy)}{\partial y}-0,\; 0-\frac{\partial(xy)}{\partial x},\; 0-0\right) = (x, -y, 0)
\]
Ce n'est pas $\vec{B}_2$. Essayons $\vec{A} = \frac{1}{2}(x^2 z - y^2 z, 0, 0)$... En fait, la propriété générale est : $\text{div}(\overrightarrow{\text{rot}}\,\vec{F}) = 0$ pour tout $\vec{F}$. On en déduit que tout champ solénoïdal peut \textit{localement} s'écrire comme un rotationnel.

\textbf{4.} Le bord de $S$ est le cercle unité $C : x^2+y^2=1$, $z=0$. Par le théorème de Stokes :
\[
\iint_S \overrightarrow{\text{rot}}\,\vec{B}_1\cdot\d\vec{S} = \oint_C \vec{B}_1\cdot\d\vec{l}
\]
Calculons d'abord $\overrightarrow{\text{rot}}\,\vec{B}_1 = -2\vec{k}$ (déjà calculé en Exercice 4). Donc :
\[
\iint_S (-2\vec{k})\cdot\d\vec{S} = -2\iint_S \d S_z
\]
Sur l'hémisphère, $\iint_S \d S_z$ = projection sur $z$ = aire du disque unité $= \pi$.

Flux $= -2\pi$.
\end{boitecorrige}

\subsection*{Exercice 10 - Laplacien en coordonnées cylindriques \niveau{\difficile}}

\begin{boiteexercice}
En coordonnées cylindriques $(r,\theta,z)$, la formule du laplacien d'un scalaire est :
\[
\Delta f = \frac{1}{r}\frac{\partial}{\partial r}\left(r\frac{\partial f}{\partial r}\right) + \frac{1}{r^2}\frac{\partial^2 f}{\partial\theta^2} + \frac{\partial^2 f}{\partial z^2}
\]
\begin{enumerate}
  \item Calculer $\Delta f$ pour $f(r) = \ln r$ et pour $f(r,\theta) = r^n\cos(n\theta)$ (harmonique cylindrique).
  \item Dans un problème de diffusion de chaleur en régime permanent en 2D ($\Delta T = 0$), avec $T$ dépendant uniquement de $r$, exprimer $T(r)$.
  \item Application : un cylindre creux de rayon interne $a=\qty{5}{cm}$ et externe $b=\qty{10}{cm}$ est maintenu à $T_a = 100\,\text{$^\circ$C}$ en $r=a$ et $T_b = 20\,\text{$^\circ$C}$ en $r=b$. Trouver $T(r)$.
\end{enumerate}
\end{boiteexercice}

\begin{boitecorrige}
\textbf{1.} Pour $f(r) = \ln r$ :
\[
\frac{\partial(\ln r)}{\partial r} = \frac{1}{r}, \quad \frac{\partial}{\partial r}\left(r\cdot\frac{1}{r}\right) = \frac{\partial(1)}{\partial r} = 0, \quad \Delta(\ln r) = \frac{1}{r}\times0 = 0
\]
$\ln r$ est harmonique (en 2D, pour $r\neq0$).

Pour $f = r^n\cos(n\theta)$ :
\[
\frac{\partial f}{\partial r} = nr^{n-1}\cos(n\theta), \quad \frac{\partial}{\partial r}(r\cdot nr^{n-1}\cos(n\theta)) = n^2r^{n-1}\cos(n\theta)
\]
\[
\frac{1}{r^2}\frac{\partial^2 f}{\partial\theta^2} = \frac{1}{r^2}(-n^2 r^n\cos(n\theta)) = -n^2r^{n-2}\cos(n\theta)
\]
\[
\Delta f = \frac{n^2r^{n-1}}{r}\cos(n\theta) - n^2 r^{n-2}\cos(n\theta) = n^2r^{n-2}\cos(n\theta) - n^2r^{n-2}\cos(n\theta) = 0
\]
Ces fonctions sont harmoniques (harmoniques cylindriques).

\textbf{2.} $\Delta T = 0$ avec $T = T(r)$ seul :
\[
\frac{1}{r}\frac{\d}{\d r}\left(r\frac{\d T}{\d r}\right) = 0 \Rightarrow r\frac{\d T}{\d r} = C_1 \Rightarrow \frac{\d T}{\d r} = \frac{C_1}{r} \Rightarrow T(r) = C_1\ln r + C_2
\]

\textbf{3.} Conditions aux limites : $T(a) = C_1\ln a + C_2 = 100$ et $T(b) = C_1\ln b + C_2 = 20$.

$C_1 \approx -115{,}4\,\mathrm{K/ln}$ et $C_2 \approx -246\,\mathrm{K}$

\[
T(r) = -115{,}4\ln r - 246 \quad (r\text{ en mètres})
\]

Vérification : $T(0{,}05) = -115{,}4\times(-2{,}996)-246 \approx 100\,^\circ\text{C}$.
\end{boitecorrige}

\subsection*{Exercice 11 - Circulation et travail \niveau{\moyen}}

\begin{boiteexercice}
Soit le champ de force $\vec{F} = (y^2, 2xy+z, y)$.
\begin{enumerate}
  \item Montrer que $\vec{F}$ est conservatif (irrotationnel).
  \item Trouver le potentiel $\varphi$ tel que $\vec{F} = \overrightarrow{\text{grad}}\,\varphi$.
  \item Calculer le travail de $\vec{F}$ le long d'un chemin quelconque de $A=(0,0,0)$ à $B=(1,2,3)$.
  \item Vérifier en paramétrant le chemin $\mathcal{C}$ : $\vec{r}(t) = (t, 2t, 3t)$ pour $t\in[0,1]$.
\end{enumerate}
\end{boiteexercice}

\begin{boitecorrige}
\textbf{1.} Calculons $\overrightarrow{\text{rot}}\,\vec{F}$ :
\begin{align*}
(\text{rot}\,\vec{F})_x &= \frac{\partial y}{\partial y} - \frac{\partial(2xy+z)}{\partial z} = 1-1=0\\
(\text{rot}\,\vec{F})_y &= \frac{\partial y^2}{\partial z} - \frac{\partial y}{\partial x} = 0-0=0\\
(\text{rot}\,\vec{F})_z &= \frac{\partial(2xy+z)}{\partial x} - \frac{\partial y^2}{\partial y} = 2y-2y=0
\end{align*}
Donc $\overrightarrow{\text{rot}}\,\vec{F} = \vec{0}$ : $\vec{F}$ est conservatif.

\textbf{2.} On cherche $\varphi$ tel que $\partial\varphi/\partial x = y^2$, $\partial\varphi/\partial y = 2xy+z$, $\partial\varphi/\partial z = y$.

De la 1re équation : $\varphi = xy^2 + g(y,z)$.
De la 2e : $\partial\varphi/\partial y = 2xy + \partial g/\partial y = 2xy+z$, donc $\partial g/\partial y = z$, soit $g = yz + h(z)$.
De la 3e : $\partial\varphi/\partial z = y + h'(z) = y$, donc $h'(z)=0$, $h = C$.
\[
\boxed{\varphi(x,y,z) = xy^2 + yz + C}
\]

\textbf{3.} Travail de $A$ à $B$ :
\[
W = \varphi(B)-\varphi(A) = (1\times4+2\times3)-0 = 4+6 = 10
\]

\textbf{4.} Sur $\mathcal{C}$ : $x=t$, $y=2t$, $z=3t$, $\d\vec{r} = (1,2,3)\d t$.
$\vec{F}(\mathcal{C}(t)) = (4t^2, 4t^2+3t, 2t)$.
\begin{align*}
W &= \int_0^1 (4t^2, 4t^2+3t, 2t)\cdot(1,2,3)\,\d t = \int_0^1(4t^2+8t^2+6t+6t)\,\d t\\
&= \int_0^1(12t^2+12t)\,\d t = \left[4t^3+6t^2\right]_0^1 = 4+6 = 10
\end{align*}
On retrouve bien $W=10$.
\end{boitecorrige}

\subsection*{Exercice 12 - Vecteur champ électrique par Gauss \niveau{\moyen}}

\begin{boiteexercice}
Une sphère creuse de rayon intérieur $a$ et extérieur $b$ porte une charge $Q$ répartie uniformément dans son volume. Utiliser le théorème de Gauss (en analogie avec le théorème de la divergence) pour calculer le champ $\vec{E}(r)$ dans les trois régions : $r<a$, $a\leq r\leq b$, $r>b$.
\end{boiteexercice}

\begin{boitecorrige}
Par symétrie sphérique, $\vec{E} = E(r)\hat{r}$. La densité volumique est $\rho = Q/V_{coque}$ avec $V_{coque} = \frac{4\pi}{3}(b^3-a^3)$.

\textbf{Pour $r < a$ (cavité) :} Aucune charge intérieure. $\oiint \vec{E}\cdot\d\vec{S} = 0$. Par symétrie, $E(r) = 0$.

\textbf{Pour $a \leq r \leq b$ (dans la coque) :} La charge intérieure est $Q_{int}(r) = \rho\cdot\frac{4\pi}{3}(r^3-a^3)$.
\[
4\pi r^2 E = \frac{Q_{int}}{\varepsilon_0} = \frac{Q(r^3-a^3)}{\varepsilon_0(b^3-a^3)}
\]
\[
\boxed{E(r) = \frac{Q}{4\pi\varepsilon_0}\cdot\frac{r^3-a^3}{r^2(b^3-a^3)}}
\]

\textbf{Pour $r > b$ (extérieur) :} Toute la charge est intérieure, $Q_{int} = Q$.
\[
\boxed{E(r) = \frac{Q}{4\pi\varepsilon_0 r^2}}
\]
La sphère se comporte comme une charge ponctuelle vue de l'extérieur.
\end{boitecorrige}


\subsection*{Exercice 13 - Théorème de Stokes : circulation et flux \niveau{\difficile}}

\begin{boiteexercice}
Soit le champ vectoriel $\vec{F}(x,y,z) = (y^2, -x, z+1)$.

\begin{enumerate}
  \item Calculer le rotationnel $\overrightarrow{\text{rot}}\,\vec{F}$.
  \item Le théorème de Stokes affirme que pour toute surface orientée $\Sigma$ de bord $\partial\Sigma$ :
  \[
  \oint_{\partial\Sigma} \vec{F}\cdot\d\vec{l} = \iint_\Sigma (\overrightarrow{\text{rot}}\,\vec{F})\cdot\d\vec{S}
  \]
  Vérifier le théorème de Stokes sur le disque $\Sigma : x^2+y^2 \leq 1$, $z=0$, orienté par $\hat{n} = \hat{z}$, et son bord $\partial\Sigma$ parcouru dans le sens trigonométrique.
  \begin{enumerate}[label=(\alph*)]
    \item Calculer le flux de $\overrightarrow{\text{rot}}\,\vec{F}$ à travers $\Sigma$.
    \item Calculer la circulation de $\vec{F}$ sur le cercle $\partial\Sigma : x = \cos t, y = \sin t, z = 0$, $t\in[0,2\pi]$.
    \item Vérifier que les deux résultats sont égaux.
  \end{enumerate}
  \item Interpréter physiquement le théorème de Stokes en termes de champ magnétique.
\end{enumerate}
\end{boiteexercice}

\begin{boitecorrige}
\textbf{1. Calcul du rotationnel.}

\[
\overrightarrow{\text{rot}}\,\vec{F} = \begin{vmatrix} \hat{x} & \hat{y} & \hat{z} \\ \partial_x & \partial_y & \partial_z \\ y^2 & -x & z+1 \end{vmatrix}
\]

\begin{align*}
(\text{rot}\,\vec{F})_x &= \frac{\partial(z+1)}{\partial y} - \frac{\partial(-x)}{\partial z} = 0 - 0 = 0 \\
(\text{rot}\,\vec{F})_y &= \frac{\partial(y^2)}{\partial z} - \frac{\partial(z+1)}{\partial x} = 0 - 0 = 0 \\
(\text{rot}\,\vec{F})_z &= \frac{\partial(-x)}{\partial x} - \frac{\partial(y^2)}{\partial y} = -1 - 2y
\end{align*}

Donc :
\[
\overrightarrow{\text{rot}}\,\vec{F} = (0, 0, -1-2y)
\]

\textbf{2.(a) Flux de $\overrightarrow{\text{rot}}\,\vec{F}$ à travers $\Sigma$ (disque $z=0$, $\hat{n} = \hat{z}$).}

\[
\iint_\Sigma (\overrightarrow{\text{rot}}\,\vec{F})\cdot\hat{z}\,\d S = \iint_{x^2+y^2\leq 1} (-1-2y)\,\d x\,\d y
\]

On passe en coordonnées polaires $x = r\cos\phi$, $y = r\sin\phi$, $\d x\,\d y = r\,\d r\,\d\phi$ :
\[
= \int_0^{2\pi}\int_0^1 (-1 - 2r\sin\phi)\,r\,\d r\,\d\phi
\]
\[
= \int_0^{2\pi}\left[\int_0^1 (-r)\,\d r - \int_0^1 2r^2\sin\phi\,\d r\right]\d\phi
\]
\[
= \int_0^{2\pi}\left[-\frac{1}{2} - \frac{2}{3}\sin\phi\right]\d\phi
\]
\[
= -\frac{1}{2}\times2\pi - \frac{2}{3}\left[-\cos\phi\right]_0^{2\pi} = -\pi - \frac{2}{3}(0) = -\pi
\]

\textbf{2.(b) Circulation de $\vec{F}$ sur $\partial\Sigma$.}

Paramétrisation : $x = \cos t$, $y = \sin t$, $z = 0$, donc $\d\vec{l} = (-\sin t, \cos t, 0)\d t$.

$\vec{F}$ sur $\partial\Sigma$ : $F_x = y^2 = \sin^2 t$, $F_y = -x = -\cos t$, $F_z = z+1 = 1$.

\begin{align*}
\oint_{\partial\Sigma} \vec{F}\cdot\d\vec{l} &= \int_0^{2\pi}\left[\sin^2 t\cdot(-\sin t) + (-\cos t)\cdot\cos t + 1\cdot 0\right]\d t \\
&= \int_0^{2\pi}\left(-\sin^3 t - \cos^2 t\right)\d t
\end{align*}

On calcule séparément :
\[
\int_0^{2\pi}\sin^3 t\,\d t = \int_0^{2\pi}\sin t(1-\cos^2 t)\,\d t = 0 \quad (\text{intégrale nulle sur une période complète})
\]
\[
\int_0^{2\pi}\cos^2 t\,\d t = \pi
\]

Donc :
\[
\oint_{\partial\Sigma} \vec{F}\cdot\d\vec{l} = 0 - \pi = -\pi
\]

\textbf{2.(c) Vérification.}

Les deux calculs donnent $-\pi$ : le théorème de Stokes est bien vérifié.
\[
\iint_\Sigma (\overrightarrow{\text{rot}}\,\vec{F})\cdot\d\vec{S} = \oint_{\partial\Sigma} \vec{F}\cdot\d\vec{l} = -\pi \quad
\]

\textbf{3. Interprétation physique.}

En électromagnétisme, la loi d'Ampère (théorème de Stokes appliqué à $\vec{B}$) stipule :
\[
\oint_C \vec{B}\cdot\d\vec{l} = \mu_0 I_{\text{enlacé}} = \iint_\Sigma \mu_0\vec{J}\cdot\d\vec{S}
\]
où $\vec{J}$ est la densité de courant. C'est exactement la forme du théorème de Stokes avec $\vec{F} = \vec{B}$ et $\overrightarrow{\text{rot}}\,\vec{B} = \mu_0\vec{J}$. La circulation du champ magnétique sur un contour fermé est égale au flux de la densité de courant à travers toute surface s'appuyant sur ce contour.
\end{boitecorrige}

\subsection*{Exercice 14 - Divergence, Stokes et coordonnées curvilignes \niveau{\difficile}}

\begin{boiteexercice}
Soit le champ $\vec{A}(r,\theta,z) = r\,\hat{r} + r\sin\theta\,\hat{z}$ en coordonnées cylindriques $(r,\theta,z)$.

\begin{enumerate}
  \item Rappeler les expressions de la divergence et du rotationnel en coordonnées cylindriques.
  \item Calculer $\text{div}\,\vec{A}$.
  \item Calculer $\overrightarrow{\text{rot}}\,\vec{A}$.
  \item Vérifier le théorème de la divergence (théorème d'Ostrogradski) sur le cylindre $\mathcal{V} : r \leq a$, $0 \leq z \leq h$, pour le champ $\vec{B}(r,\theta,z) = r\,\hat{r}$.
  \begin{enumerate}[label=(\alph*)]
    \item Calculer $\text{div}\,\vec{B}$ en coordonnées cylindriques.
    \item Calculer le volume intégral $\iiint_\mathcal{V} \text{div}\,\vec{B}\,\d V$.
    \item Calculer le flux de $\vec{B}$ à travers la surface fermée $\partial\mathcal{V}$ (face latérale + deux disques).
    \item Vérifier l'égalité.
  \end{enumerate}
\end{enumerate}
\end{boiteexercice}

\begin{boitecorrige}
\textbf{1. Opérateurs en coordonnées cylindriques.}

Pour un champ $\vec{F} = F_r\hat{r} + F_\theta\hat{\theta} + F_z\hat{z}$ :
\[
\text{div}\,\vec{F} = \frac{1}{r}\frac{\partial(rF_r)}{\partial r} + \frac{1}{r}\frac{\partial F_\theta}{\partial\theta} + \frac{\partial F_z}{\partial z}
\]
\[
(\overrightarrow{\text{rot}}\,\vec{F})_r = \frac{1}{r}\frac{\partial F_z}{\partial\theta} - \frac{\partial F_\theta}{\partial z}, \quad
(\overrightarrow{\text{rot}}\,\vec{F})_\theta = \frac{\partial F_r}{\partial z} - \frac{\partial F_z}{\partial r}, \quad
(\overrightarrow{\text{rot}}\,\vec{F})_z = \frac{1}{r}\frac{\partial(rF_\theta)}{\partial r} - \frac{1}{r}\frac{\partial F_r}{\partial\theta}
\]

\textbf{2. Divergence de $\vec{A} = r\,\hat{r} + r\sin\theta\,\hat{z}$.}

Ici $A_r = r$, $A_\theta = 0$, $A_z = r\sin\theta$.

\[
\text{div}\,\vec{A} = \frac{1}{r}\frac{\partial(r\cdot r)}{\partial r} + 0 + \frac{\partial(r\sin\theta)}{\partial z}
= \frac{1}{r}\cdot 2r + 0 = 2
\]

\textbf{3. Rotationnel de $\vec{A}$.}

\[
(\overrightarrow{\text{rot}}\,\vec{A})_r = \frac{1}{r}\frac{\partial(r\sin\theta)}{\partial\theta} - \frac{\partial(0)}{\partial z} = \frac{1}{r}\cdot r\cos\theta = \cos\theta
\]
\[
(\overrightarrow{\text{rot}}\,\vec{A})_\theta = \frac{\partial r}{\partial z} - \frac{\partial(r\sin\theta)}{\partial r} = 0 - \sin\theta = -\sin\theta
\]
\[
(\overrightarrow{\text{rot}}\,\vec{A})_z = \frac{1}{r}\frac{\partial(r\cdot 0)}{\partial r} - \frac{1}{r}\frac{\partial r}{\partial\theta} = 0 - 0 = 0
\]

Donc : $\overrightarrow{\text{rot}}\,\vec{A} = \cos\theta\,\hat{r} - \sin\theta\,\hat{\theta}$.

\textbf{4. Théorème d'Ostrogradski pour $\vec{B} = r\,\hat{r}$.}

\textbf{4.(a)} $B_r = r$, $B_\theta = B_z = 0$ :
\[
\text{div}\,\vec{B} = \frac{1}{r}\frac{\partial(r^2)}{\partial r} = \frac{2r}{r} = 2
\]

\textbf{4.(b)} Volume intégral ($\d V = r\,\d r\,\d\theta\,\d z$ en cyl.) :
\[
\iiint_\mathcal{V} 2\,\d V = 2\times\text{Vol}(\mathcal{V}) = 2\times\pi a^2 h
\]

\textbf{4.(c)} Flux à travers $\partial\mathcal{V}$.

\textit{Face latérale} ($r = a$, $0\leq\theta\leq2\pi$, $0\leq z\leq h$, $\d\vec{S} = a\,\d\theta\,\d z\,\hat{r}$) :
\[
\Phi_{lat} = \int_0^h\int_0^{2\pi} a\cdot a\,\d\theta\,\d z = a^2\cdot 2\pi\cdot h = 2\pi a^2 h
\]

\textit{Disques supérieur et inférieur} ($\d\vec{S} = \pm\hat{z}\,r\,\d r\,\d\theta$) : comme $\vec{B} = r\,\hat{r}$ n'a pas de composante $\hat{z}$, le flux à travers les disques est $\Phi_{disques} = 0$.

Total : $\Phi_{total} = 2\pi a^2 h + 0 = 2\pi a^2 h$.

\textbf{4.(d) Vérification :}
\[
\iiint_\mathcal{V} \text{div}\,\vec{B}\,\d V = 2\pi a^2 h = \oiint_{\partial\mathcal{V}} \vec{B}\cdot\d\vec{S} \quad
\]
Le théorème d'Ostrogradski est bien vérifié.
\end{boitecorrige}

\ifdefined\inMainDoc\else\end{document}\fi
